\documentclass[11pt,a4paper]{article}
\usepackage[margin=2.2cm]{geometry}
\usepackage[T1]{fontenc}
\usepackage{lmodern,amsmath,amssymb,booktabs,longtable,array,enumitem}
\usepackage[colorlinks=true,urlcolor=blue,linkcolor=blue]{hyperref}
\expandafter\def\expandafter\UrlBreaks\expandafter{\UrlBreaks\do\_\do\-}
\emergencystretch=3em
\newcommand{\A}{\textsf{\textbf{(a)}}}
\newcommand{\B}{\textsf{\textbf{(b)}}}
\newcommand{\C}{\textsf{\textbf{(c)}}}
\newcommand{\D}{\textsf{\textbf{(d)}}}
\newcommand{\file}[1]{\nolinkurl{#1}}
\newcommand{\ours}{\file{paper/scripts/review_draft1_checks.py}}
\newcolumntype{L}[1]{>{\raggedright\arraybackslash}p{#1}}
\title{Independent review of the first scattering-paper draft}
\author{Review for PI Gustav Delius}
\date{4 November 2026}
\begin{document}
\maketitle

\noindent\textbf{Recommendation: revise before circulation.} \D\
The core classical phase survives the independent checks \C.
The nine appendix items do not invalidate the computed classical scattering
phase, but two have consequences beyond a local typo: the bad cubic sign
reappears in a normal-ordering energy claim, and the signed defect-root
dictionary misidentifies the physical defect for negative orientation \C.
There are \textbf{14 required corrections}, including coordinate and charge
definitions and qualifications of stability and fusion. The corrected good
classical energy was not used incorrectly in the inspected one-loop
calculations \C. Passing algebraic checks does not settle the physical
one-loop determinant or quantum frame \D.

\paragraph{Scope and provenance.}
This review reads \file{AGENTS.md}, \file{README.md}, the full draft,
the five requested source write-ups/map, the four named one-loop/charge
write-ups, relevant downstream notes, and the cited script implementations.
Only the three requested deliverables are changed; no git operation is used.
Labels retain their repository meanings. Review findings, numerical
counterexamples and source/code observations are \C; recommendations and
unestablished physical extensions are \D. All new mathematical calculations
below are checked in \ours; the run groups are named A1--A9 and B.
Source-script evidence is distinguished from an independent proof.

The original bibliography's reading labels belong to earlier project agents.
For this review, the Corrigan--Zambon full-text sections 2--5 were read online,
especially (3.22), (3.37), (3.43)--(3.44), and (5.5).
The Bogdan--Charkina full text was consulted at the classical energy,
breather and action formulas; not every page was reread. The
Charlier--Lenells water-wave abstract/entry was independently verified;
the large HTML full text could not be retrieved. Other literature claims
are assessed against the supplied literature map and source write-ups,
not certified by a fresh full-text literature review. No new reference is
invented. The precise reading limitations matter to corrections R12--R13.

\section*{A. Triage of all nine appendix items}
\small
\begin{longtable}{@{}L{1.4cm}L{3.7cm}L{3.9cm}L{6.0cm}@{}}
\toprule Item & Correct version \C & Evidence & Downstream effect \C\\\midrule
\endfirsthead
\toprule Item & Correct version \C & Evidence & Downstream effect \C\\\midrule
\endhead
1. Bad cubic &
$\mathcal L_b=\frac12\phi_t^2-\frac12\phi_x^2-\phi_x^3+\frac12\phi_{xx}^2$.
$E_b=\frac2{15}k^3(5+4k^2)$. &
A1: exact Euler--Lagrange calculation; exact pulse integrals and 65-digit direct quadrature. &
Classical phase, $M,\Pi,P$ and the computed energy use the correct data.
However, \file{phase2/oneloop.tex} and the docstring of
\file{oneloop_soliton_energy.py} give the opposite sign for the normal-ordering
boundary energy. Their displayed numerical tables use scheme term zero;
those tables are unaffected. See A.1 below.\\\addlinespace
2. Good energy &
$E_g=\frac2{15}k^3(5-4k^2)$; $dE_g=v\,dP$.
The printed one-loop energy exceeds this by $k^3(5-2k^2)/15$. &
A2: symbolic integration, direct field-energy quadrature for both directions;
actual downstream function bodies evaluated independently. &
None found in the specified one-loop renormalisation, charge or phase
calculations. \file{audit_charges_action.py:ecl} explicitly uses the correct
energy. The incorrect displayed energy and the adjacent asserted derivative
identity must both be corrected. This does not clear the separate physical
one-loop audits.\\\addlinespace
3. Defect root &
At fixed CZ fields and orientation: $s=\sigma_0$, $q=-s/(2\sqrt3)$,
$\alpha_{CZ}=\eta/\sqrt3$, $\beta_{CZ}=2\eta/\sqrt3$.
The spectrum may separately use $\eta\sigma_{vac}$. &
A3: CZ constraints and nonlinear sewing on a one-soliton for both $\eta$ and
both directions. A signed-root substitution at fixed defect fails. &
Wrong physical root labels propagate into the text of
\file{oneloop_defect_phase.tex} and Route E of \file{report.tex}.
The two \file{oneloop_defect_*} calculation scripts use the unsigned root.
Their functions and formal chord identities are unaffected; interpretation
as functions of the same CZ $\sigma$ for $\eta=-1$ is affected.\\\addlinespace
4. Transmission &
$z=((v+s-\sqrt3 k)/(v+s+\sqrt3 k))^\eta
=((p_b-q)/(p_a-q))^\eta$.
The printed middle expression is $-1/z$. &
A4: rational identity for both orientations; comparison with actual
\file{z_defect_mode}; nonlinear sewing. &
None found in executable transmission, delay, phase or fusion formulas.
They use the rational/root expression, not the incorrect middle expression.
It is a real sign and reciprocal error, not a harmless logarithm branch.\\\addlinespace
5. KdV ratio &
The \emph{full} derivative ratio is $0.564246498175\ldots$ at ratio $1.7$;
the cross-only ratio is $-0.435753501825\ldots$. &
A5: independent high-precision small-amplitude derivative with scale
$10^{-12}$; asymptotic expression below. &
The master-report prose is wrong, and the final explanatory print line of
\file{consistency_checks.py} has a minus-sign typo. Its computed numerical
ratio and the paper's precise ratio are correct. No phase value changes.\\\addlinespace
6. Script names &
\file{smatrix_poles_bootstrap.py}, \file{toda_quantum_corrections.py} are the
existing names. &
A6: filesystem inventory; old basenames absent, replacements present under
\file{phase2/scripts}. &
Documentation/reproducibility only. The two omitted S-matrix checks exist.
All explicitly cited script paths in paper Sections 1--6 exist.\\\addlinespace
7. Root notation &
Ordered $\mu_1=-p_b$, $\mu_2=-p_a$;
$\lambda=-\ell/(12\sqrt3)$. Bad roots instead obey
$\mu^3+\mu/4=i\ell/(12\sqrt3)$. &
A7 and paper conventions check 4; roots substituted in the two cubics and
cross-ratio checked in both real forms. &
None for this notation change. It requires an ordered dictionary, not just
an unidentified statement ``$\mu=-p$''. Additional bad transcription
formulas in Phase 2g are noted below; the paper does not copy them.\\\addlinespace
8. Toda width &
$\kappa_H=2\kappa_T$, and
$\kappa_H\sim\sqrt{12}\,k_b$, $\kappa_T\sim\sqrt3\,\beta$ at small amplitude. &
A8: exact factor two in the exponential and
$\sinh\beta/\beta\to1$; source lattice formulas and scaling checks. &
None. These are asymptotic identifications; $k_b=\sinh\beta$ must not be
replaced by $\beta$ at finite amplitude.\\\addlinespace
9. Phase number &
\file{defect_semiclassics.pdf} is Phase 2c. &
A9: title, README and source inventory. &
None. Phase 2d in the one-loop defect narrative is a citation typo.\\
\bottomrule
\end{longtable}
\normalsize

\subsection*{A.1 Energy signs and the downstream trace}
For $u=k^2\operatorname{sech}^2(kx/2)/2$ and $\pi=-vu$, direct integration gives
\C\ (A1--A2)
\[
 \int u^2dx=\frac23k^3,\qquad \int u^3dx=\frac4{15}k^5,
 \qquad \int u_x^2dx=\frac2{15}k^5.
\]
Inserting these in the two Hamiltonians gives the energy formulas in the
table. With the wrong bad cubic sign, the energy on the same positive pulse
would be $2k^3/3$, differing from the correct value by $-8k^5/15$; moreover
that pulse is not a solution of that wrong equation \C\ (A1).
The wrong good energy has derivative residual
\C\ (A2)
\[
 \frac{dE_{\rm printed}}{dk}-v\frac{dP}{dk}
 =k^2\left(1-\frac23k^2\right)\ne0\quad(0<k<1).
\]
This form avoids division by $P'$ at the cusp.

The bad Hamiltonian contains $+u^3$. Under the usual Wick replacement
$u^3\mapsto:u^3:=u^3-3Cu$, its energy shift is
\C\ (A1)
\[
 \Delta E_{\rm Wick}=-3CM
 =-\frac{3\hbar}{2\pi}M
 =-\frac{3\hbar}{\pi}\sinh\beta,
 \qquad C=\frac{\hbar}{2\pi}
\]
for the previously selected stable-band vacuum. The positive shift asserted
in \file{phase2/oneloop.tex}, the paragraph ``Decision (PI)'' and the energy
table caption, is the sign appropriate to the \emph{good} Hamiltonian.
The same positive claim occurs at lines 15--16 of
\file{phase2/scripts/oneloop_soliton_energy.py}. Its runtime code only computes
$C$ and the zero-scheme integral; it does not add either boundary term.
\file{oneloop_breather.py} also uses the zero-scheme values.
Thus this is propagation into a claimed renormalised energy, not evidence
that the numerical zero-scheme tables or classical phase were computed with
the wrong Lagrangian \C. No new scheme is selected by this review.

For the good-energy typo, the dependency trace is:
\begin{itemize}[leftmargin=*]
\item \file{oneloop_good_phonon_scattering.py} builds modes from the PDE and
Hirota data, with no classical energy input. \C
\item \file{oneloop_good_twosoliton_phase.py} computes the continuum integrals,
\file{Delta_P1_closed}, \file{Delta_E1_ren_closed}, and
\file{delta_renorm_good}. The last uses derivatives of the classical phase
and momentum. Independent checks reproduce it in both collision channels.
\file{oneloop_good_limits_cusp_resonance.py} imports those quantities. \C
\item \file{charge_and_frame_audit.py} computes
$L_1=v\Delta P_{mode}-\Delta E_1$ from mode sums and differentiates it.
The new script evaluates these exact function bodies and verifies their
agreement with the preceding mode charges. No $E_{\rm printed}$ is used. \C
\item \file{audit_charges_action.py:ecl} is
$2\sin^3\alpha(1+4\cos^2\alpha)/15$, the correct classical energy.
Its finite-dimensional charge argument and the corresponding write-up do
not inherit the typo. \C
\end{itemize}
The remaining charge/frame and determinant limitations in the independent
audits are separate issues; correcting $E_{\rm printed}$ does not resolve
them \D.

\subsection*{A.2 The defect convention is settled, not an open choice}
CZ's full-text equations impose $\beta_{CZ}=2\alpha_{CZ}$ and
$\beta_{CZ}^2-\alpha_{CZ}^2=1$ in the stated units \A.
The later pair $\alpha_{CZ}=1/(2\sqrt3)$,
$\beta_{CZ}=\eta/(2\sqrt3)$ violates both constraints \C\ (A3).
It is not an alternative parametrisation of that same action.
The source's $s$ in (5.5) is $\sigma_0$ when $c=1$.
\href{https://arxiv.org/html/2306.03772v1}{Corrigan--Zambon, full-text sections 3 and 5}.

For example, $\alpha=0.4$, $\epsilon=+1$, $\sigma_0=2.4$, $\eta=-1$
gives $z=1.50971083399\ldots$. Replacing the root by
$-\eta\sigma_0/(2\sqrt3)$ gives $z=0.373565030471\ldots$;
the maximum residual of the four sewing equations at $t=0.27,x=0$ is
$3.3312537\ldots$, while the correct factor passes below $10^{-52}$ \C\ (A3).
Thus transmission, delay, absorption-root identification and root-labelled
phase/charge values change if the later dictionary is applied literally.
The spectrum legitimately depends on $\eta\sigma_{vac}$, but that combination
must not replace $\sigma_{vac}$ in the scattering root without a simultaneous
parity transformation of the bulk data \C.

Both \file{oneloop_defect_scattering_and_stability.py} and
\file{oneloop_defect_phase_and_fusion.py} retain $s=-2\sqrt3q$ and the
correct rational/root factor; their executable functions do not use the
incorrect middle expression. The later master report copies the bad
dictionary in its Route E prose. The mirrored-triple loop in
\file{audit_frames_defects.py} uses an extra $\eta$ in its local variable
\file{sigmas}; it still counts one stable root, since both orientations and
reflected triples are sampled, but does not establish equality of the two
physical root dictionaries \C. Its determinant checks use the signed
spectral parameter legitimately.

Two nearby Phase 2g transcription errors should not be imported when repairing
the dictionary: its displayed cubic has $\ell/(4\sqrt3)$ in place of
$\ell/(12\sqrt3)$, and its trigonometric expressions for $p_a,p_b$ do not
equal their preceding Cartesian expressions. The Cartesian roots used by
the code and by the paper are correct \C\ (A7/B root substitutions).

\subsection*{A.3 Minor-item precision}
The independent KdV limit is \C\ (A5)
\[
 \frac{\partial_{\beta_1}\delta}{2\beta_1\log A}
 \longrightarrow 1+\frac{2\beta_2}{\beta_1\log A_{KdV}},
 \qquad A_{KdV}=\left(\frac{\beta_1-\beta_2}{\beta_1+\beta_2}\right)^2.
\]
The extra $4\beta_2$ is the cross-term contribution to the \emph{phase
derivative}, namely $-\Pi_1'M_2$, not the dimensionally different spatial
shift itself. The paper mostly repairs the earlier prose, but its small
amplitude paragraph should use this precise name. Script inventories are
filesystem observations, not mathematical tests; A6 and A9 test their stated
existence/title assertions rather than fabricating symbolic checks.

\section*{B. Required corrections to the paper}
Locations refer to source labels and subsection names, so they remain usable
if another agent changes line numbers. Each item states a review finding
\C\ and a recommended fix \D. The list contains exactly 14 items.
\begin{enumerate}[label=R\arabic*.,leftmargin=*,itemsep=8pt]
\item \textbf{Proposition \file{prop:EP}, one-soliton coordinate.}
Equation \file{eq:1sol} defines $x_0$ as the additive exponent, not the
centre. At $t=0$, $X=-x_0/k$ and the pullback is
$\Omega=dP\wedge dX=-P_k\,dk\wedge dx_0/k$.
Replace the printed $\pm P_k\,dk\wedge dx_0$ by this formula, or explicitly
rename the physical centre $X$. Direct quadrature gives $-0.78$ at $k=0.6$,
where $P_k=0.468$ (B). The degeneracy at the cusp survives.
The later phase scripts use physical positions and do not inherit this
normalisation error.

\item \textbf{Section 2, Gelfand--Dickey charges.}
The displayed definition $H_n=\int\operatorname{Res}L^{n/3}dx$ is divergent
on the nonzero Lax background. Already
$\operatorname{Res}L^{1/3}\to-1/12$ (B).
Use the vacuum-subtracted residue, as \file{charges.py:H} actually does.
The quoted finite $H_1,H_2,H_4,H_5$ and their $E,P$ combinations then agree
with the code; label this definition explicitly before using them.

\item \textbf{Section 2, paragraph after \file{eq:shifts}, and Section 4
small amplitudes.}
The KdV expression printed as a position shift is $\log A_{KdV}$.
For the stated ordering, write
$\Delta x_1\sim-2k_1^{-1}\log|(k_1-k_2)/(k_1+k_2)|$ and
$\Delta x_2\sim+2k_2^{-1}\log|(k_1-k_2)/(k_1+k_2)|$.
Also call $-\Pi_1'M_2\to4\beta_2$ a contribution to the phase derivative,
as distinguished in A.3. The exact shifts and phase are unaffected.

\item \textbf{Section 4, Toda set-up.}
Restore the direction sign in $dE/dP_n=s b\sinh\kappa/\kappa$.
The unsigned expression is false for left-movers (B).
It is inherited from \file{toda_route.tex}; the executable Toda phase
functions use the signed velocity, so the head-on scaling test survives.

\item \textbf{Section 5, defect symplectic structure.}
With the stated canonical one-form and $\Omega=d\Theta$, the contribution
is $(g_u-f_v)du\wedge dv$, not $(f_v-g_u)du\wedge dv$.
Change that sign or explain a changed convention. Its restriction to a fixed
vacuum sector is still zero, so no soliton--defect cross term is introduced.
Do not infer absence of localized defect energy or momentum from this
vanishing two-form; the paper already correctly calls $M_D,\Pi_D$ effective
phase coefficients.

\item \textbf{Abstract, summary, \file{prop:stab}, and the ensuing consequences.}
Qualify the stability theorem by its evidence: the cited finite root search
is not a proof of absence of all unstable modes. State the sampled result
as \C\ and the unsupplied completeness assertion as \D, or provide an
admissible-root/completeness proof. Separate the marginal case
$\eta\sigma=-1$ from strict stability. Restrict ``one stable vacuum per
triple'' to three distinct real roots. For example $s=2$ has vacua
$2,-1,-1$; $|s|>2$ has only one real vacuum (B).
The unrestricted statement ``both excited vacua are unstable'' is likewise
inapplicable when those vacua are nonreal. These qualifications are already
recognised in \file{audit_frames_defects.tex}.

\item \textbf{Proposition \file{prop:fusion}, absolute constant.}
Replace ``with no free constant'' by ``for the displayed zero-constant
representatives''. A shift $\delta_D^{(q)}\mapsto\delta_D^{(q)}+C(q)$ leaves
every classical delay invariant and adds $C(q')-C(q)$ to the fusion
identity. The cited script sets those constants to zero; it does not
determine an absolute physical phase. Preserve the algebraic identity as
\C. Give the orientation sign explicitly, e.g.
$\rho=\epsilon_3\eta\operatorname{sgn}(q-q')$ for the ordered chord
convention used by the scripts, instead of leaving it undetermined.

\item \textbf{``Labels and verification'', verification footnotes and appendix.}
Replace the assertion that every \C\ formula already has a matching asserted
test by a precise coverage statement. In particular:
\file{charges.py} does not check the $x_0$ symplectic formula;
\file{defect_two_soliton_symplectic.py}, check 5, prints limits and a
\emph{defect} cusp but does not assert the paper's two-body square-root
claim; and the Toda two-body reflection amplitude is explicitly ``not
separately checked'' in the appendix. The present B checks supply a
Bessel-wave test for that last amplitude. The Hirota exact-arithmetic test
covers selected rational amplitudes in every direction pattern for
$N\le4$, not a symbolic proof at all amplitudes; one separate profile
assertion in that script is disabled by \file{or True}. State the actual
scope, and do not count printed diagnostics as assertions. The legacy
\file{phase0/checks/convention_check.py} also lacks the required final
success marker, although its sign assertions pass.

\item \textbf{Appendix ``Inconsistencies found'', all nine items.}
Replace unresolved guesses with table A's verdicts and downstream trace.
In particular disclose the bad normal-ordering sign, the verified CZ
parameter constraints, and the negative-orientation physical mismatch;
remove ``could not check'' and ``whether it propagated ... was not checked''.
Retain the distinction between text errors and unaffected executable
calculations. Record the minor final-print sign error in
\file{consistency_checks.py} and qualify the Toda width map as asymptotic.

\item \textbf{Sections 2 and 4, source dictionaries.}
The paper's own good/bad dictionary and the CZ factor-nine dictionary are
correct, but the imported BC and CL formulas also need explicit field maps.
For BC use $\varphi_{BC}=-\phi_{bad}$ and
$w_{BC}=-u_{bad}$; this leaves the canonical action/energy normalisation
unchanged. For the CL bad PDE use $u_{CL}=3u_{bad}$; for the BD field
normalisation used in Phase 1, $u_{BD}=u_{good}/2$.
Identify the bad roots called $p$ in \file{eq:crossratio} with the preceding
$\mu_{a,b}$. Otherwise the ``both equations'' cross-ratio uses an undefined
bad-equation $p$. These are notation/attribution repairs, not changes to
the verified continuation of $\delta$.

\item \textbf{Section 1, quantum prior work, affine rank.}
For $a_n^{(1)}$ the underlying finite algebra is
$\mathfrak{sl}_{n+1}$, not $\mathfrak{sl}_n$.
Write $U_q(\widehat{\mathfrak{sl}}_{n+1})$, or use the consistent alternative
notation $a_{N-1}^{(1)}$ with $\mathfrak{sl}_N$.
The Section 6 $a_2/\mathfrak{sl}_3$ formulas themselves are consistent.

\item \textbf{Section 6, obstructions; Section 1, no-boost inference.}
The observation that the displayed affine-Toda and Boussinesq constants
have different mass/dispersion terms does not prove that \emph{no} limit
of fixed-species affine Toda can produce Boussinesq. No cited script proves
that universal exclusion. Replace it by the supported statement that no
such dictionary/limit has been constructed here (\D\ for a conjectured
obstruction). Likewise absence of a boost is not alone a proof against
every possible reparametrised difference form. State the checked
nonseparability in the chosen rapidities, or supply the stronger proof.
Describe the QISM issue as an obstruction to the ordinary ultralocal
construction, not to every possible lattice quantisation.

\item \textbf{Literature labels in Sections 1, 2 and 4.}
Separate the literature's explicit classical action calculation (\A)
from its semiclassical Bohr--Sommerfeld prescription (\B), and keep the
project's string identity conditional on that prescription (\C).
The source expressly assumes large quantum number; it does not prove an
exact quantum spectrum for every $N$. Specify hypotheses/function spaces
when paraphrasing the CL existence theorem and the abstract-only
Bona--Sachs claim. Do not present the literature-search absence result as
a theorem of nonexistence. Verify the bibliography entries explicitly
marked ``not machine-verified'' before making the blanket assertion that
all references are verified. Historical reading-status labels may stay if
clearly attributed to the earlier agents, as the draft's last paragraph does.

\item \textbf{Summary item 6 and claims-table equation tags.}
Split the combined ``good \& bad'' summary into what was actually checked:
the canonical non-ultralocal bracket is displayed for good, the
interaction-block factorisation is checked in both forms, and the cited
product-form exclusion script tests bad overtaking. Either supply the
additional good exclusion check or narrow that tag. The hierarchy-charge
row should also distinguish the good $H_n$ evaluation from the two-form
root dictionary. Toda itself is a lattice model; ``bad'' denotes its
small-amplitude continuum target, not an identity of finite-amplitude
theories.
\end{enumerate}

\section*{Displayed-equation and claim-label ledger (Sections 1--6)}
Every displayed environment in the requested sections is covered below;
inline formulas with a material defect are included in the corrections.
``Pass'' means source agreement plus the specified check, not a claim of
independent proof for every continuum or literature assertion. All audit
assessments in this ledger are \C.

\small
\begin{longtable}{@{}L{5.1cm}L{4.9cm}L{5.6cm}@{}}
\toprule Location / displayed formula & Source and check & Verdict / label scope\\\midrule
\endfirsthead
\toprule Location / displayed formula & Source and check & Verdict / label scope\\\midrule
\endhead
Section 1 (no displayed equation environments in its body) &
Literature map; Phase 1 introduction; master-report scope. &
Prior-work (a)/(b) mostly follows the map; R11--R14 qualify rank,
inference and equation attribution. Abstract phase agrees with Section 4.\\
\file{eq:good}, \file{eq:bad}, \file{eq:Lag}, \file{eq:Omega0}, \file{eq:charges} &
Phase 1 equations/canonical definitions; paper check 1; independent A1--A2. &
Pass. The correct bad Lagrangian is the paper's improvement over its source.
Potential EL equation is differentiated once to obtain the field PDE.\\
\file{eq:1sol}, \file{eq:rapidity}, \file{eq:solcharges}, \file{eq:dP} &
Phase 1; Hirota/charges scripts; paper checks 2,3,9; independent A2. &
Pass; interpret $dE/dP$ by continuation at the cusp. Inline symplectic
coordinate needs R1.\\
\file{eq:roots}, \file{eq:cubic}; inline bad cubic and $H_n$ &
Phase 1 charges, Phase 2c transmission and bad ansatz;
paper check 4; independent B. &
Root equations pass. The $H_n$ values pass after R2 subtraction.
The $\mu=-p$ relation reverses ordering.\\
\file{eq:Nsol} &
Phase 1/Hirota. Exact rational-amplitude checks through $N=4$. &
Agrees with source. General $N$ is literature (a), sampled project check (c).
Do not call a finite rational grid a general proof; R8.\\
\file{eq:Agood}, \file{eq:Abad}, \file{eq:crossratio} &
Phase 1 interaction regimes; consistency check 5; independent B. &
Pass in both forms and orientations, with bad roots named explicitly (R10).
Exclude coincident solitons and singular denominators when using ordinary
two-body scattering.\\
\file{eq:shifts}; inline breather constant &
Phase 1 peak tracking and breather scan; paper conventions. &
Exact shifts pass for soliton 1 initially left and $v_1>v_2$.
Good breather bound is only for the stated Hirota zero-background class.
The following KdV sentence needs R3.\\
Dictionary table; unnumbered continuation display &
Phase 1 Lax script; paper check 5. &
Pass as an analytic identity of formulas, not a map between real solutions.
CZ exclusion is only within their nonlinear type-I ansatz.\\
\file{eq:Omega}, \file{eq:sympid} &
Master report; symplectic scattering and two-soliton defect scripts. &
Pass as asymptotic pullbacks, with fixed global potential constant and
well-separated solitons. Exponentially small terms are explicitly excluded.\\
\file{eq:darboux}, \file{eq:darbouxshifts} &
Master report; semiclassical/consistency scripts; independent B. &
Pass locally off $P'=0$. In--out difference uses the ordering change.
Normalising the opposite potential endpoint gives the same differences.\\
\file{eq:delta} &
Master report; independent high-precision gradients, both equations and
all four sign combinations. &
Pass as a classical generating function in the specified frame.
An exact quantum S-matrix is not established by this identity.\\
Inline properties of \file{prop:props} &
Paper checks 6--7 and existing defect script check 5. &
Antisymmetry and regular-branch crossing pass. Resonance finiteness follows
from the vanishing spectral difference; cusp citation does not assert its
two-body local expansion (R8). Logarithm sheet and regular region matter.\\
\file{eq:Ibr}, \file{eq:string} &
BC action formula; consistency check 4; paper check 8. &
Classical action identity passes. Quantisation is semiclassical and
conditional (R13); exact pole positions remain (d).\\
Toda unnumbered $M,\Pi,E$ display; inline $P_n,A_T$ &
\file{toda_route.tex}, lattice-soliton and label scripts. &
Pass; derivative velocity needs R4 sign. Vacuum energy and pressure-work
subtractions are already stated in the paper.\\
Toda inline $R,\Theta$ and classical limit &
Route D and quantum BA; independent Bessel/Schrodinger test B. &
Reflection amplitude passes with incoming $e^{-iqx/\hbar}$ and
$\Delta p=2q$. Choose an unwrapped logarithm for the classical limit.
The BA method remains labelled (b).\\
\file{prop:todaBA}, \file{eq:deltaT}, unnumbered scaling ratio &
Classical BA, dressed-phase, frame, label and scaling scripts. &
Matches the source and numerical checks. Mixed-derivative primes refer
to the respective $\kappa_j$; not arbitrary physical-momentum derivatives.
Numerical agreement is not a complete analytic convergence proof.\\
\file{eq:zCZ}, \file{eq:zroot} &
CZ (5.5), Phase 2c transmission; independent A3/A4 sewing. &
Pass with $s=\sigma_0$. Require finite positive $z$ for ordinary transmission;
the endpoints are absorption limits.\\
Inline $\Omega_D$ and \file{eq:dD} &
Phase 2c transmission checks 4 and 6; exact one-form calculation. &
Phase and gradient pass. R5 corrects the unrestricted two-form sign;
its vacuum restriction remains zero. $s$-dependent constants are unfixed.\\
\file{prop:2sD}; vacuum root/absorption identities &
Phase 2c one- and two-soliton scripts. &
Pass for the regular sampled configurations; algebraic factorisation is
consistent with the preserved $A_{12}$. Excited-state interpretation needs
the real-vacuum restrictions in R6.\\
\file{eq:fusion} &
\file{defect_smatrix_good_continuation.py}, check 4. &
Pass for zero-constant representatives; R7. The script determines the
relative orientation sign at sampled triangles, not a quantum bootstrap.\\
Section 6 unnumbered Poisson brackets &
\file{routes_BC_checks.py}; independent smeared bracket B. &
Pass. $\{r,r\}=0$ follows by integration by parts for decaying test functions.
The nonzero derivative-delta bracket establishes non-ultralocality.\\
\file{eq:fact}; unnumbered $G(y)$; unnumbered good factorisation &
Route B/C checks and interaction regimes; independent B for Toda/$G$ blocks. &
Pass. The good expression is the trigonometric continuation with its own
$S,D$. No all-limits affine-Toda exclusion follows (R12).\\
Section 6 product-form exclusion (inline mixed derivative) &
\file{smatrix_rejected_forms.py}. &
The nonzero mixed derivative excludes $f(D)+g(S)$ in the tested bad
rapidities. Narrow the summary tag or add the good check (R14).\\
\bottomrule
\end{longtable}
\normalsize

\paragraph{Label audit beyond the displays.}
The explicit project formulas in Sections 2--6 are correctly classified
as (c), subject to the corrections and test coverage just stated.
The mutual-statistics interpretation, universality, cusp uniformisation,
bound-state poles, resonance interpretation and exact S-matrix ideas are
appropriately (d). The Toda/GHD and affine-Toda quantum proposals are
appropriately (b). The missing-spin statement is algebraic and needs no
unread textbook as its sole warrant. The opening assertion that the quantum
Boussinesq theory ``is'' the $W_3$ charges should be read as describing that
literature construction, not uniqueness of all possible quantisations.
The broad stability, fusion-constant and affine-limit claims are the main
places where the paper goes beyond its evidence. The first two already
appear overstrong in older sources; copying them is not an independent
justification \C.

\paragraph{Phase 2i exclusion.}
The draft contains no Phase 2i Wiener--Hopf factors, edge coefficients,
two-loop formula, corrected kernel or new physical synthesis in Sections
1--6, and no such result is imported into the remainder of the draft \C.
Its Section 8 placeholder names \file{toda_wh_and_kp_synthesis.pdf} solely
as work awaiting audit. The quantum placeholders otherwise trace to the
earlier audits already cited in the master report. This is a status mention,
not use of Phase 2i results. It can be removed if the PI intends a complete
ban on even mentioning the filename; this review imports none of its claims.

\section*{Optional improvements}
\begin{itemize}[leftmargin=*]
\item Separate field, velocity and Toda sign notation: $u_b,V,\epsilon_T$
avoid the repeated use of $v$ and $s$. \D
\item Give a short analytic derivation of the phase-gradient identity and
of the real-root count. The numerical evidence is strong, but an explicit
proof is more useful to a reader than a long collection of precisions. \D
\item Define scattering order once ($v_1>v_2$) and distinguish
$\delta(1,2)$ as an antisymmetric label function from the phase assigned
to a specific physical incoming ordering. \D
\item Add a compact branch/endpoint convention for logarithms, the cusp,
coincident solitons and breather boundary limits. For bad solitons say
$E$ increases with $|P|$, rather than saying $E(P)$ is monotone without
specifying a direction branch. \D
\item Keep the quantum placeholders short. In particular, the identification
of a retained linear phonon band with the full nonlinear CL invariant
subspace is not independently proved by a finite phase integral; retain
the heuristic qualification from \file{oneloop.tex}. \D
\end{itemize}

\section*{Checks, reproducibility and failures recorded}
\paragraph{Independent checks.}
Run \file{python3 -B paper/scripts/review_draft1_checks.py}.
The script writes no files, uses exact SymPy and 65-digit mpmath, asserts
the new identities and explicit counterexamples, inventories citations,
and ends with \file{ALL CHECKS PASSED}. Actual source function definitions
are loaded via the Python AST without executing their top-level code or
rewriting their result files. This makes the downstream comparison
independent of saved JSON outputs. The optional \file{--run-sources} reruns
the read-only scripts cited before Section 7, omitting the figure writer.

\paragraph{Original check requested by the PI.}
\file{paper/scripts/paper_conventions_check.py} was run unchanged and printed
\file{ALL CHECKS PASSED}. Its checks establish the paper's energy and
convention fixes, not the downstream dependency claims or every assertion
in Sections 1--6.

\paragraph{Source-script coverage.}
The cited source scripts were read, and the read-only verification scripts
for Sections 1--6 were rerun with bytecode writes disabled. The figure
generator was inspected but not rerun, because it would overwrite an
existing figure. Numerical root-search completeness and references not
freshly read are not certified by these runs. The full Phase 2 one-loop and
audit suites were not indiscriminately rerun: several write existing result
files, contrary to the task constraint. The relevant actual functions were
instead checked as described above. All 24 source-script processes exited
successfully; 23 printed \file{ALL CHECKS PASSED}. The legacy
\file{phase0/checks/convention_check.py} asserted its two symbolic residuals
and printed the correct signs, but lacks the required success marker.
Together with the separately run paper script, this covers the 25
verification scripts among the 26 cited paths (the remaining path is the
figure generator). The independent review script also printed
\file{ALL CHECKS PASSED}. \C

\paragraph{Failures recorded (including repaired review checks).}
\begin{enumerate}[leftmargin=*]
\item The report's bad cubic sign fails the target Euler--Lagrange equation
and positive-soliton energy. The associated positive Wick shift is also
wrong for the corrected bad Hamiltonian. These are retained as explicit
counterexamples, not repaired in the source files. \C
\item The printed good energy fails direct quadrature and the velocity
identity. The signed-root substitution at $\eta=-1$ fails nonlinear sewing;
the middle transmission expression fails equality to $z$. The independent
checks assert these nonzero residuals. \C
\item The paper's phase-coordinate symplectic coefficient and unsigned
left-moving Toda velocity fail direct checks. The unqualified strict
stable-root count fails at a repeated-root threshold; the unsubtracted
first hierarchy charge has a nonzero background density. \C
\item In the first review-script run, the marginal-root test reused a
SymPy symbol declared positive; the solver therefore excluded the desired
root $-1$. A new real symbol fixes that test without changing its expected
root or tolerance. \C
\item An exploratory repository search used an invalid regular expression;
it was replaced by narrower/fixed-string searches. One guessed script
basename did not exist; the actual script inventory was then used.
An initial full-text URL requested a nonexistent CZ version; version 1
was retrieved successfully. BC HTML and one PDF screenshot retrieval
failed; its PDF text was available. The large CL HTML and the attempted
Bona--Sachs DOI fetch were unavailable. These access failures limit fresh
literature verification as stated above. \C
\item Source verification contains limitations that must remain visible:
the disabled \file{or True} assertion in the Hirota script, printed-only
limit/cusp diagnostics, and the numerical stability search. No source
assertion or tolerance was weakened by this review. \C
\item The first PDF build failed because \file{xurl.sty} is unavailable;
the document was changed to use the installed URL package's line-breaking
hooks. A patch attempt with a stale context line was rejected without
modifying files and was reapplied with the correct context. The legacy
Phase 0 script's missing success marker is a reporting failure, not a
failed mathematical residual. \C
\end{enumerate}

\paragraph{Verified external entries used in this review.}
Corrigan and Zambon, \emph{Adding integrable defects to the Boussinesq equation},
\href{https://arxiv.org/abs/2306.03772}{arXiv:2306.03772}; full-text relevant
sections read. Bogdan and Charkina, \emph{Quasiclassical quantization of the
Boussinesq breather emerging from the kink localized mode},
\href{https://arxiv.org/abs/2502.14695}{arXiv:2502.14695}; full-text relevant
formula passages consulted. Charlier and Lenells, \emph{On Boussinesq's
equation for water waves},
\href{https://arxiv.org/abs/2204.02365}{arXiv:2204.02365}; abstract/entry
verified in this review, not a fresh full-text reading.
\end{document}
